\documentclass[12pt,a4paper]{article}
\usepackage[UTF8]{ctex}
\usepackage{amsmath,amssymb,amsthm}
\usepackage{geometry}
\usepackage{graphicx}
\usepackage{hyperref}
\usepackage{bm}

\geometry{margin=2.5cm}

\title{反对称矩阵恒等式$[\bm{a}][\bm{b}] - [\bm{b}][\bm{a}] = [[\bm{a}]\bm{b}]$的详细推导}
\author{Modern Robotics}
\date{}

\begin{document}

\maketitle

\section{引言}

在机器人动力学中，反对称矩阵（skew-symmetric matrix）的运算经常出现。本文将详细推导一个重要的恒等式：

\begin{equation}
[\bm{a}][\bm{b}] - [\bm{b}][\bm{a}] = [[\bm{a}]\bm{b}]
\label{eq:main_identity}
\end{equation}

其中$\bm{a}, \bm{b} \in \mathbb{R}^3$是三维向量，$[\bm{a}]$表示向量$\bm{a}$的反对称矩阵表示。

\section{反对称矩阵的定义}

对于向量$\bm{a} = (a_x, a_y, a_z)^T \in \mathbb{R}^3$，其反对称矩阵定义为：

\begin{equation}
[\bm{a}] = \begin{bmatrix}
0 & -a_z & a_y \\
a_z & 0 & -a_x \\
-a_y & a_x & 0
\end{bmatrix}
\label{eq:skew_definition}
\end{equation}

这个矩阵具有以下性质：
\begin{itemize}
\item 反对称性：$[\bm{a}]^T = -[\bm{a}]$
\item 叉积表示：$[\bm{a}]\bm{b} = \bm{a} \times \bm{b}$，对于任意$\bm{b} \in \mathbb{R}^3$
\end{itemize}

\section{恒等式的详细推导}

\subsection{方法1：直接矩阵乘法}

\subsubsection{计算$[\bm{a}][\bm{b}]$}

设$\bm{a} = (a_x, a_y, a_z)^T$，$\bm{b} = (b_x, b_y, b_z)^T$，则：

\begin{align}
[\bm{a}][\bm{b}] &= \begin{bmatrix}
0 & -a_z & a_y \\
a_z & 0 & -a_x \\
-a_y & a_x & 0
\end{bmatrix}
\begin{bmatrix}
0 & -b_z & b_y \\
b_z & 0 & -b_x \\
-b_y & b_x & 0
\end{bmatrix} \\
&= \begin{bmatrix}
-a_z b_z - a_y b_y & a_y b_x & a_z b_x \\
a_z b_y & -a_x b_x - a_z b_z & a_x b_y \\
-a_y b_z & a_x b_z & -a_x b_x - a_y b_y
\end{bmatrix}
\label{eq:ab_product}
\end{align}

详细计算第一行第一列元素：
\begin{align}
([\bm{a}][\bm{b}])_{11} &= 0 \cdot 0 + (-a_z) \cdot b_z + a_y \cdot (-b_y) \\
&= -a_z b_z - a_y b_y
\end{align}

第一行第二列元素：
\begin{align}
([\bm{a}][\bm{b}])_{12} &= 0 \cdot (-b_z) + (-a_z) \cdot 0 + a_y \cdot b_x \\
&= a_y b_x
\end{align}

其他元素类似计算。

\subsubsection{计算$[\bm{b}][\bm{a}]$}

类似地：

\begin{align}
[\bm{b}][\bm{a}] &= \begin{bmatrix}
0 & -b_z & b_y \\
b_z & 0 & -b_x \\
-b_y & b_x & 0
\end{bmatrix}
\begin{bmatrix}
0 & -a_z & a_y \\
a_z & 0 & -a_x \\
-a_y & a_x & 0
\end{bmatrix} \\
&= \begin{bmatrix}
-b_z a_z - b_y a_y & b_y a_x & b_z a_x \\
b_z a_y & -b_x a_x - b_z a_z & b_x a_y \\
-b_y a_z & b_x a_z & -b_x a_x - b_y a_y
\end{bmatrix}
\label{eq:ba_product}
\end{align}

\subsubsection{计算$[\bm{a}][\bm{b}] - [\bm{b}][\bm{a}]$}

将式(\ref{eq:ab_product})和式(\ref{eq:ba_product})相减：

\begin{align}
[\bm{a}][\bm{b}] - [\bm{b}][\bm{a}] &= \begin{bmatrix}
0 & a_y b_x - b_y a_x & a_z b_x - b_z a_x \\
a_z b_y - b_z a_y & 0 & a_x b_y - b_x a_y \\
-a_y b_z + b_y a_z & a_x b_z - b_x a_z & 0
\end{bmatrix}
\end{align}

注意到：
\begin{align}
a_y b_x - b_y a_x &= -(\bm{a} \times \bm{b})_z = -([\bm{a}]\bm{b})_z \\
a_z b_x - b_z a_x &= (\bm{a} \times \bm{b})_y = ([\bm{a}]\bm{b})_y \\
a_z b_y - b_z a_y &= -(\bm{a} \times \bm{b})_x = -([\bm{a}]\bm{b})_x \\
a_x b_y - b_x a_y &= (\bm{a} \times \bm{b})_z = ([\bm{a}]\bm{b})_z \\
-a_y b_z + b_y a_z &= (\bm{a} \times \bm{b})_x = ([\bm{a}]\bm{b})_x \\
a_x b_z - b_x a_z &= -(\bm{a} \times \bm{b})_y = -([\bm{a}]\bm{b})_y
\end{align}

因此：
\begin{equation}
[\bm{a}][\bm{b}] - [\bm{b}][\bm{a}] = \begin{bmatrix}
0 & -([\bm{a}]\bm{b})_z & ([\bm{a}]\bm{b})_y \\
([\bm{a}]\bm{b})_z & 0 & -([\bm{a}]\bm{b})_x \\
-([\bm{a}]\bm{b})_y & ([\bm{a}]\bm{b})_x & 0
\end{bmatrix} = [[\bm{a}]\bm{b}]
\label{eq:result}
\end{equation}

\subsection{方法2：使用叉积性质}

利用反对称矩阵与叉积的关系：$[\bm{a}]\bm{b} = \bm{a} \times \bm{b}$。

对于任意向量$\bm{c} \in \mathbb{R}^3$，考虑：

\begin{align}
([\bm{a}][\bm{b}] - [\bm{b}][\bm{a}])\bm{c} &= [\bm{a}]([\bm{b}]\bm{c}) - [\bm{b}]([\bm{a}]\bm{c}) \\
&= [\bm{a}](\bm{b} \times \bm{c}) - [\bm{b}](\bm{a} \times \bm{c}) \\
&= \bm{a} \times (\bm{b} \times \bm{c}) - \bm{b} \times (\bm{a} \times \bm{c})
\end{align}

使用向量三重积的恒等式：
\begin{align}
\bm{a} \times (\bm{b} \times \bm{c}) &= (\bm{a} \cdot \bm{c})\bm{b} - (\bm{a} \cdot \bm{b})\bm{c} \\
\bm{b} \times (\bm{a} \times \bm{c}) &= (\bm{b} \cdot \bm{c})\bm{a} - (\bm{b} \cdot \bm{a})\bm{c}
\end{align}

因此：
\begin{align}
([\bm{a}][\bm{b}] - [\bm{b}][\bm{a}])\bm{c} &= (\bm{a} \cdot \bm{c})\bm{b} - (\bm{a} \cdot \bm{b})\bm{c} - (\bm{b} \cdot \bm{c})\bm{a} + (\bm{b} \cdot \bm{a})\bm{c} \\
&= (\bm{a} \cdot \bm{c})\bm{b} - (\bm{b} \cdot \bm{c})\bm{a}
\end{align}

使用另一个向量三重积恒等式：
\begin{equation}
(\bm{a} \cdot \bm{c})\bm{b} - (\bm{b} \cdot \bm{c})\bm{a} = (\bm{a} \times \bm{b}) \times \bm{c}
\end{equation}

因此：
\begin{align}
([\bm{a}][\bm{b}] - [\bm{b}][\bm{a}])\bm{c} &= (\bm{a} \times \bm{b}) \times \bm{c} \\
&= [\bm{a} \times \bm{b}]\bm{c} \\
&= [[\bm{a}]\bm{b}]\bm{c}
\end{align}

由于这对任意$\bm{c}$都成立，我们得到：
\begin{equation}
[\bm{a}][\bm{b}] - [\bm{b}][\bm{a}] = [[\bm{a}]\bm{b}]
\end{equation}

\subsection{方法3：使用Lie括号}

在Lie代数$\mathfrak{so}(3)$中，Lie括号定义为：
\begin{equation}
[[\bm{a}], [\bm{b}]] = [\bm{a}][\bm{b}] - [\bm{b}][\bm{a}]
\end{equation}

同时，Lie括号满足：
\begin{equation}
[[\bm{a}], [\bm{b}]] = [[\bm{a}]\bm{b}]
\end{equation}

这是因为$\mathfrak{so}(3)$与$\mathbb{R}^3$（配备叉积）同构。

\section{应用示例}

\subsection{在机器人动力学中的应用}

在牛顿-欧拉逆向动力学中，这个恒等式用于推导加速度递推公式。

考虑伴随变换的时间导数：
\begin{equation}
\frac{d}{dt}[\text{Ad}_{T}] = [\text{ad}_V] \text{Ad}_T
\end{equation}

其中$V$是twist，$\text{ad}_V$的推导中使用了这个恒等式。

具体地，当计算$\frac{d}{dt}[[p]R]$时：
\begin{align}
\frac{d}{dt}[[p]R] &= [\dot{p}]R + [p]\dot{R} \\
&= [[\omega]p + v]R + [p][\omega]R \\
&= [\omega][p]R - [p][\omega]R + [v]R + [p][\omega]R \\
&= [\omega][p]R + [v]R
\end{align}

这里使用了恒等式$[\omega][p] - [p][\omega] = [[\omega]p]$。

\section{总结}

本文详细推导了反对称矩阵的恒等式：
\begin{equation}
[\bm{a}][\bm{b}] - [\bm{b}][\bm{a}] = [[\bm{a}]\bm{b}]
\end{equation}

这个恒等式在机器人动力学中非常重要，特别是在推导加速度递推公式和计算伴随变换的时间导数时。

\end{document}

